\documentclass[12pt,a4paper]{article}

\usepackage[utf8]{inputenc}
\usepackage[T1]{fontenc}
\usepackage[spanish]{babel}
\usepackage{amsmath, amssymb}
\usepackage{geometry}
\usepackage{booktabs}
\usepackage{array}
\usepackage{xcolor}
\usepackage{hyperref}
\usepackage{enumitem}
\usepackage{fancyhdr}
\usepackage{graphicx}
\usepackage{listings}
\usepackage{tcolorbox}
\usepackage{multirow}
\usepackage{caption}

\geometry{top=2.5cm, bottom=2.5cm, left=3cm, right=3cm}

% ── Colores ──────────────────────────────────────────────────────────────────
\definecolor{uddblue}{RGB}{0, 75, 135}
\definecolor{lightgray}{RGB}{245, 245, 245}
\definecolor{codegray}{RGB}{100, 100, 100}
\definecolor{codegreen}{RGB}{0, 120, 0}
\definecolor{codeblue}{RGB}{0, 0, 180}
\definecolor{codered}{RGB}{160, 0, 0}

% ── Estilo de código Python ───────────────────────────────────────────────────
\lstset{
    language     = Python,
    backgroundcolor = \color{lightgray},
    basicstyle   = \small\ttfamily,
    keywordstyle = \color{codeblue}\bfseries,
    commentstyle = \color{codegreen}\itshape,
    stringstyle  = \color{codered},
    numberstyle  = \tiny\color{codegray},
    numbers      = left,
    stepnumber   = 1,
    breaklines   = true,
    frame        = single,
    rulecolor    = \color{gray!40},
    tabsize      = 4,
    showstringspaces = false,
    literate =
        {á}{{\'{a}}}1 {é}{{\'{e}}}1 {í}{{\'{i}}}1
        {ó}{{\'{o}}}1 {ú}{{\'{u}}}1 {ñ}{{\~{n}}}1
        {Á}{{\'{A}}}1 {É}{{\'{E}}}1 {Í}{{\'{I}}}1
        {Ó}{{\'{O}}}1 {Ú}{{\'{U}}}1 {Ñ}{{\~{N}}}1
}

% ── Cajas de ayuda ────────────────────────────────────────────────────────────
\tcbuselibrary{skins, breakable}
\newtcolorbox{ayuda}[1]{
    colback   = lightgray,
    colframe  = uddblue,
    fonttitle = \bfseries\small,
    title     = {#1},
    breakable,
    left      = 6pt,
    right     = 6pt,
    top       = 4pt,
    bottom    = 4pt
}

% ── Encabezado y pie ─────────────────────────────────────────────────────────
\pagestyle{fancy}
\fancyhf{}
\fancyhead[L]{\small\color{gray} IIP314W\,--\,Optimización Aplicada a Negocios}
\fancyhead[R]{\small\color{gray} Tarea 4\,--\,2026-T1}
\fancyfoot[C]{\small\thepage}
\renewcommand{\headrulewidth}{0.4pt}

\hypersetup{colorlinks=true, linkcolor=uddblue, urlcolor=uddblue}

% ── Espaciado entre párrafos ──────────────────────────────────────────────────
\setlength{\parskip}{0.5em}
\setlength{\parindent}{0pt}

% =============================================================================
\begin{document}
% =============================================================================

% ── Título ───────────────────────────────────────────────────────────────────
\begin{center}
    {\Large\bfseries\color{uddblue}
     (IIP314W) Optimización Aplicada a Negocios}\\[0.4em]
    {\large Universidad del Desarrollo}\\[0.8em]
    \noindent\rule{\linewidth}{1.5pt}\\[0.5em]
    {\LARGE\bfseries Tarea 4}\\[0.3em]
    {\large Análisis de Sensibilidad y Precios Sombra}\\[0.5em]
    \noindent\rule{\linewidth}{1.5pt}\\[0.8em]
    \begin{tabular}{@{}ll@{}}
        \textbf{Profesor:} & Ing.\ Rodrigo Trigo Vilches \\[0.3em]
        \textbf{Ayudante:} & Lic.\ Vicente Ramírez Almonacid \\[0.3em]
        \textbf{Período:}  & Año 2026\,--\,Trimestre 1
    \end{tabular}
\end{center}

\vspace{1em}
\noindent\rule{\linewidth}{0.4pt}
\vspace{0.5em}

% =============================================================================
\section*{Contexto: FrutaSur Procesados S.A.}
% =============================================================================

FrutaSur Procesados S.A.\ es una empresa mediana de la Región del Maule dedicada
a la transformación de frutas de temporada en productos de valor agregado.
Abastece a cadenas de supermercados y distribuidores a nivel nacional.

La empresa fabrica tres líneas de productos:

\begin{itemize}
    \item \textbf{Jugo Natural (JN):} zumo embotellado sin conservantes elaborado
          directamente desde fruta fresca.
          Margen de contribución: \$120 por caja.
    \item \textbf{Pulpa Congelada (PC):} pulpa procesada y sellada al vacío,
          destinada a panaderías y restaurantes.
          Margen de contribución: \$95 por caja.
    \item \textbf{Mermelada Artesanal (MA):} elaborada con cocción lenta y
          azúcar de caña.
          Margen de contribución: \$180 por caja.
\end{itemize}

La planta opera de lunes a sábado con cuatro recursos productivos cuya
disponibilidad semanal es fija:

\begin{center}
\renewcommand{\arraystretch}{1.25}
\begin{tabular}{lcccc}
\toprule
\textbf{Recurso} &
\textbf{JN (por caja)} &
\textbf{PC (por caja)} &
\textbf{MA (por caja)} &
\textbf{Disponible semanal} \\
\midrule
Fruta fresca (kg)    & 8     & 12   & 6    & 960  \\
Azúcar (kg)          & 0{,}5 & 0    & 3    & 180  \\
Horas de procesado   & 2     & 1{,}5 & 4   & 280  \\
Horas de envasado    & 1     & 2    & 1{,}5 & 200 \\
\bottomrule
\end{tabular}
\captionof{table}{Requerimientos de recursos por unidad producida.
Fuente: elaboración propia.}
\end{center}

El objetivo de la empresa es \textbf{maximizar el margen de contribución
semanal total}, sujeto a las disponibilidades de recursos indicadas.

\noindent\rule{\linewidth}{0.4pt}

% =============================================================================
\section*{Parte A \quad Formulación del Modelo \hfill\normalfont\normalsize(20 pts)}
% =============================================================================

Formule el \textbf{modelo de programación lineal} correspondiente al problema
de FrutaSur. Su formulación debe incluir:

\begin{enumerate}[label=\alph*), leftmargin=*, topsep=2pt]
    \item \textbf{Variables de decisión:} defínalas con precisión, indicando
          unidades y dominio.
    \item \textbf{Función objetivo:} expresión matemática del margen total a
          maximizar.
    \item \textbf{Restricciones:} una por cada recurso, con su nombre o
          etiqueta al costado. Incluya restricciones de no negatividad.
\end{enumerate}

\noindent\textit{Nota: el modelo debe quedar completamente expresado en
notación matemática estándar, no en código.}

\noindent\rule{\linewidth}{0.4pt}

% =============================================================================
\section*{Parte B \quad Resolución con Gurobi \hfill\normalfont\normalsize(20 pts)}
% =============================================================================

Implemente el modelo de la Parte~A en \texttt{gurobipy} y resuélvalo.
Responda:

\begin{enumerate}[label=\alph*), leftmargin=*, topsep=2pt]
    \item Reporte la \textbf{solución óptima}: cantidad de cajas de cada
          producto a producir semanalmente.
    \item Reporte el \textbf{valor óptimo} de la función objetivo.
    \item ¿Qué restricciones están \textbf{activas} (holgura\,=\,0) en la
          solución óptima? ¿Cuáles tienen holgura positiva y cuánto vale?
    \item Interprete en \textbf{lenguaje de negocios}: ¿qué significa que
          una restricción tenga holgura cero? ¿Y holgura positiva?
\end{enumerate}

\begin{ayuda}{Ayuda -- Holguras en Gurobi}
\vspace{2pt}
Luego de \texttt{model.optimize()}, acceda a las holguras de las restricciones:
\begin{lstlisting}
for c in model.getConstrs():
    print(f"{c.ConstrName:20s}  slack = {c.Slack:.4f}")
\end{lstlisting}
Para que el nombre sea informativo, asegúrese de nombrarlo al crear la
restricción:
\begin{lstlisting}
model.addConstr(8*x_JN + 12*x_PC + 6*x_MA <= 960,
                name="fruta_fresca")
\end{lstlisting}
\end{ayuda}

\noindent\rule{\linewidth}{0.4pt}

% =============================================================================
\section*{Parte C \quad Análisis de Sensibilidad \hfill\normalfont\normalsize(40 pts)}
% =============================================================================

% ── C.1 ──────────────────────────────────────────────────────────────────────
\subsection*{C.1 \quad Precios Sombra \hfill\normalfont\normalsize(15 pts)}

\begin{enumerate}[label=\alph*), leftmargin=*, topsep=2pt]
    \item Obtenga el \textbf{precio sombra} $\pi_i$ de cada restricción.
    \item Para cada restricción activa, interprete $\pi_i$ en lenguaje de
          negocios: ¿cuánto aumentaría el margen semanal si se dispusiera
          de una unidad adicional de ese recurso?
    \item Identifique el recurso con el \textbf{mayor precio sombra}:
          ¿en cuál convendría invertir prioritariamente para aumentar la
          rentabilidad semanal?
    \item ¿Por qué las restricciones con holgura positiva tienen precio
          sombra igual a cero? Explique la intuición económica.
\end{enumerate}

\begin{ayuda}{Ayuda -- Precios sombra en Gurobi}
\begin{lstlisting}
print(f"{'Restriccion':<20}  {'Pi (precio sombra)':>20}")
print("-" * 44)
for c in model.getConstrs():
    print(f"{c.ConstrName:<20}  {c.Pi:>20.4f}")
\end{lstlisting}
\texttt{c.Pi} entrega el valor dual de la restricción, es decir, el cambio
en el valor óptimo por unidad adicional del RHS\@ (precio sombra~$\pi_i$).
Solo tiene sentido consultarlo cuando el modelo está en estado
\texttt{GRB.OPTIMAL}.
\end{ayuda}

\vspace{0.5em}
\newpage
% ── C.2 ──────────────────────────────────────────────────────────────────────
\subsection*{C.2 \quad Rango de la Función Objetivo \hfill\normalfont\normalsize(15 pts)}

\begin{enumerate}[label=\alph*), leftmargin=*, topsep=2pt]
    \item Obtenga el \textbf{rango de variación del coeficiente de la FO}
          para cada variable: ¿cuánto puede subir o bajar el margen de
          un producto sin que cambie el plan óptimo de producción?
    \item Presente los resultados en una tabla con columnas:
          \emph{Producto | Margen actual | Límite inferior | Límite superior}.
    \item Analice el producto con el rango \textbf{más estrecho}:
          ¿qué tan sensible es el plan óptimo a variaciones en su margen?
          ¿Qué implicancia tiene para la gerencia de ventas?
    \item Analice el producto con el rango \textbf{más amplio}:
          ¿qué significa en términos prácticos que su coeficiente pueda
          variar ampliamente sin alterar la solución?
\end{enumerate}

\begin{ayuda}{Ayuda -- Ranging de la función objetivo en Gurobi}
\begin{lstlisting}
print(f"{'Variable':<10}  {'Obj':>8}  {'SAObjLow':>12}  {'SAObjUp':>12}")
print("-" * 46)
for v in model.getVars():
    print(f"{v.VarName:<10}  {v.Obj:>8.2f}  "
          f"{v.SAObjLow:>12.4f}  {v.SAObjUp:>12.4f}")
\end{lstlisting}
\texttt{v.SAObjLow} y \texttt{v.SAObjUp} son los límites del intervalo en
que puede variar el coeficiente de \texttt{v} en la FO \emph{manteniendo
la misma base óptima} (y por tanto el mismo plan de producción).
\end{ayuda}

\vspace{0.5em}
\newpage
% ── C.3 ──────────────────────────────────────────────────────────────────────
\subsection*{C.3 \quad Rango del Lado Derecho (RHS) \hfill\normalfont\normalsize(10 pts)}

\begin{enumerate}[label=\alph*), leftmargin=*, topsep=2pt]
    \item Obtenga el \textbf{rango de variación del RHS} para cada
          restricción: ¿en qué intervalo puede variar la disponibilidad
          de cada recurso para que los precios sombra de C.1 sigan siendo
          válidos?
    \item Presente los resultados en una tabla con columnas:
          \emph{Restricción | RHS actual | Límite inferior | Límite superior}.
    \item La empresa tiene la opción de adquirir fruta fresca adicional en
          el mercado spot al precio de \$4 por kg. Según el análisis de
          sensibilidad: ¿conviene comprarla? ¿Hasta cuántos kg adicionales
          sería rentable adquirir? Justifique su respuesta con los valores
          obtenidos.
\end{enumerate}

\begin{ayuda}{Ayuda -- Ranging del RHS en Gurobi}
\begin{lstlisting}
print(f"{'Restriccion':<20}  {'RHS':>8}  {'SARHSLow':>12}  {'SARHSUp':>12}")
print("-" * 56)
for c in model.getConstrs():
    print(f"{c.ConstrName:<20}  {c.RHS:>8.2f}  "
          f"{c.SARHSLow:>12.4f}  {c.SARHSUp:>12.4f}")
\end{lstlisting}
\texttt{c.SARHSLow} y \texttt{c.SARHSUp} entregan el intervalo en que puede
variar el RHS de la restricción~\texttt{c} sin que cambie la \emph{base
óptima} (y por tanto sin que cambie el precio sombra). Fuera de este rango,
la base cambia y los precios sombra dejan de ser válidos.
\end{ayuda}

\noindent\rule{\linewidth}{0.4pt}
\newpage
% =============================================================================
\section*{Parte D \quad Análisis Paramétrico \hfill\normalfont\normalsize(20 pts)}
% =============================================================================

% ── D.1 ──────────────────────────────────────────────────────────────────────
\subsection*{D.1 \quad Variación del margen de la Mermelada Artesanal}

El equipo comercial ha informado que el margen de la MA podría fluctuar
entre \$80 y \$280 por caja según las condiciones del mercado.

\begin{enumerate}[label=\alph*), leftmargin=*, topsep=2pt]
    \item Resuelva el modelo para \textbf{al menos 40 valores} del margen
          de la MA en el rango $[80,\,280]$. Registre el valor óptimo de
          la FO y la solución $(x_{JN}^*,\,x_{PC}^*,\,x_{MA}^*)$ para cada
          valor.
    \item Grafique el \textbf{valor óptimo} (eje $y$) en función del margen
          de la MA (eje $x$). Identifique y anote en el gráfico los
          \textbf{puntos de quiebre} (valores donde cambia el plan óptimo).
    \item En un segundo gráfico, trace las \textbf{cantidades óptimas de
          cada producto} (tres curvas en el mismo eje) en función del margen
          de la MA.
    \item ¿En qué valor del margen la MA deja de producirse (o comienza
          a producirse)? Compare este valor con el rango obtenido en la
          Parte~C.2. ¿Coinciden? ¿Por qué?
\end{enumerate}

\vspace{0.5em}

% ── D.2 ──────────────────────────────────────────────────────────────────────
\subsection*{D.2 \quad Variación de la disponibilidad de fruta fresca}

\begin{enumerate}[label=\alph*), leftmargin=*, topsep=2pt]
    \item Resuelva el modelo para \textbf{al menos 40 valores} de la
          disponibilidad de fruta fresca en el rango $[400,\,1600]$\,kg.
          Registre el valor óptimo de la FO para cada valor.
    \item Grafique el \textbf{valor óptimo} en función de la disponibilidad
          de fruta. Identifique y marque los \textbf{puntos de quiebre}
          (cambios de pendiente).
    \item ¿Cómo se relaciona la \textbf{pendiente} del gráfico, en cada
          tramo, con el precio sombra de la restricción de fruta calculado
          en C.1? ¿Qué ocurre con la pendiente cuando la restricción pasa
          de activa a inactiva?
    \item Verifique que los rangos del RHS obtenidos en C.3 se
          \textbf{corresponden con los puntos de quiebre} observados en
          el gráfico. ¿Coinciden?
\end{enumerate}
\newpage
\begin{ayuda}{Ayuda -- Análisis paramétrico en Gurobi}
Para variar un parámetro y resolver el modelo múltiples veces:
\begin{lstlisting}
import numpy as np
import matplotlib.pyplot as plt

# Silenciar output de Gurobi antes del loop
model.Params.OutputFlag = 0

# --- D.1: variar coeficiente de la FO ---
margenes_MA = np.linspace(80, 280, 50)
valores_opt  = []
plan_JN, plan_PC, plan_MA = [], [], []

for c_MA in margenes_MA:
    x_MA.Obj = c_MA          # actualizar coeficiente
    model.update()
    model.optimize()
    if model.Status == GRB.OPTIMAL:
        valores_opt.append(model.ObjVal)
        plan_JN.append(x_JN.X)
        plan_PC.append(x_PC.X)
        plan_MA.append(x_MA.X)

# Restaurar coeficiente original
x_MA.Obj = 180
model.update()

# --- D.2: variar RHS de fruta fresca ---
constr_fruta = model.getConstrByName("fruta_fresca")
disponibilidades = np.linspace(400, 1600, 50)
valores_opt_fruta = []

for b in disponibilidades:
    constr_fruta.RHS = b     # actualizar RHS
    model.update()
    model.optimize()
    if model.Status == GRB.OPTIMAL:
        valores_opt_fruta.append(model.ObjVal)

# Restaurar RHS original
constr_fruta.RHS = 960
model.update()
\end{lstlisting}
Recuerde restaurar los parámetros originales al final de cada loop para no
arrastrar cambios a la siguiente parte.
\end{ayuda}

\noindent\rule{\linewidth}{0.4pt}

% =============================================================================
\section*{Entregables}
% =============================================================================

\subsection*{1.\quad Informe (PDF)}

Estructura mínima requerida:

\begin{enumerate}[leftmargin=*, topsep=2pt]
    \item \textbf{Portada:} nombre del curso, Tarea~4, integrantes, fecha.
    \item \textbf{Índice.}
    \item \textbf{Introducción:} descripción del problema y objetivos del
          análisis. 
    \item \textbf{Marco teórico:} Conceptos importantes que se deben definir
          para entender el problema y desarrollo.
    \item \textbf{Desarrollo:} Partes A, B, C y D con sus respectivas
          preguntas respondidas.
    \item \textbf{Resultados:} Outputs, valores numericos, tablas, figuras, etc.
    \item \textbf{Conclusiones:} Principales observaciones, explicaciones e interpretaciones.
    \item \textbf{Referencias} (si aplica).
\end{enumerate}

\textbf{Formalidades} (se descuentan puntos por incumplimiento):

\begin{itemize}[topsep=2pt]
    \item Lenguaje técnico y formal, sin errores ortográficos ni gramaticales.
    \item Texto justificado, fuente tamaño 12.
    \item Cada sección comienza en una nueva página.
    \item Cada figura, gráfico o tabla debe tener \textbf{título y fuente}
          (elaboración propia si corresponde).
\end{itemize}

\subsection*{2.\quad Cuaderno Jupyter (.ipynb)}

\begin{itemize}[topsep=2pt]
    \item Implementación del modelo en \texttt{gurobipy} (Parte~B).
    \item Extracción e impresión de holguras, precios sombra y rangos
          (Parte~C).
    \item Código del análisis paramétrico con los gráficos solicitados
          (Parte~D).
    \item Comentarios explicativos en cada sección del código.
    \item Código limpio, ordenado y bien estructurado.
\end{itemize}

\medskip
\noindent\textbf{Formato de entrega:} comprima el informe y el notebook en
un archivo \texttt{.zip} o \texttt{.rar} con el nombre
\texttt{Tarea4\_GrupoN}. \textbf{No entregue los archivos por separado.}

\noindent\rule{\linewidth}{0.4pt}

% =============================================================================
\section*{Plazos y Evaluación}
% =============================================================================

\begin{itemize}[topsep=2pt]
    \item \textbf{Plazo de entrega:} Viernes 15 de mayo a las 23:59.
    \item \textbf{Grupos:} 2 a 3 integrantes (se recomienda 3).
    \item \textbf{Inscripción:} anótense en los grupos de Canvas
          (Personas~$\to$ Grupos~$\to$ Tarea~4\,--\,Grupo~n), incluso
          si trabajan solos.
\end{itemize}

\vspace{0.5em}

\begin{center}
\renewcommand{\arraystretch}{1.2}
\begin{tabular}{lc}
\toprule
\textbf{Criterio} & \textbf{Ponderación} \\
\midrule
Parte A: Formulación del modelo                      & 20\,\% \\
Parte B: Resolución e interpretación                 & 20\,\% \\
Parte C: Análisis de sensibilidad                    & 40\,\% \\
Parte D: Análisis paramétrico y gráficos             & 20\,\% \\
\midrule
\textit{Calidad del informe y profundidad del análisis} &
\textit{transversal} \\
\bottomrule
\end{tabular}
\captionof{table}{Criterios de evaluación. Fuente: elaboración propia.}
\end{center}

\vspace{0.3em}
\noindent\textit{La calidad del informe, la profundidad del análisis y la
correcta interpretación de los resultados en lenguaje de negocios impactan
de forma transversal en todas las partes.}

\noindent\rule{\linewidth}{0.4pt}

% =============================================================================
\section*{Notas Importantes}
% =============================================================================

\begin{itemize}[topsep=2pt]
    \item Los atributos de sensibilidad de Gurobi (\texttt{.Pi},
          \texttt{.Slack}, \texttt{.SAObjLow}, etc.) solo están disponibles
          \textbf{luego de resolver} y solo si el estado es
          \texttt{GRB.OPTIMAL}.
    \item En el análisis paramétrico, \textbf{silenciee el output} de
          Gurobi (\texttt{model.Params.OutputFlag = 0}) antes del loop
          para evitar saturar la salida.
    \item Al finalizar cada loop paramétrico, \textbf{restaure el valor
          original} del parámetro variado para no arrastrar cambios entre
          partes.
    \item En el informe y en el código deben figurar los \textbf{nombres
          de todos los integrantes del grupo}.
    \item \textbf{Deben usar \texttt{gurobipy}} para la implementación.
\end{itemize}

\vspace{1em}
\begin{center}
\textit{Consultas al correo
\href{mailto:vramireza@udd.cl}{\texttt{vramireza@udd.cl}} o al
\texttt{+56\,9\,9912\,1814}
(respondo siempre que puedo, dentro de horarios razonables).}
\end{center}

% =============================================================================
\end{document}
% =============================================================================
